\documentclass{amsart}

\usepackage{amsmath,amssymb}

\newtheorem{theorem}{Theorem}
\newtheorem{lemma}{Lemma}
\newtheorem{proposition}{Proposition}

\DeclareMathOperator{\Cov}{Cov}

\begin{document}

\begin{lemma}
For \(s>1\), the Dirichlet eta function
\[
\eta(s):=(1-2^{1-s})\zeta(s)
\]
is nondecreasing.
\end{lemma}

\begin{proof}
For \(s>1\),
\[
\eta(s)
=
\frac{1}{\Gamma(s)}
\int_0^\infty \frac{x^{s-1}}{e^x+1}\,dx.
\]
Let \(X_s\) have the gamma distribution with density
\[
\frac{x^{s-1}e^{-x}}{\Gamma(s)}
\qquad (x>0),
\]
and put
\[
h(x):=\frac{1}{1+e^{-x}}.
\]
Then
\[
\eta(s)=\mathbb E[h(X_s)].
\]
Differentiating under the integral sign gives
\[
\eta'(s)
=
\Cov\bigl(h(X_s),\log X_s\bigr).
\]
Since both \(h(x)\) and \(\log x\) are increasing functions of \(x\),
\[
\Cov\bigl(h(X_s),\log X_s\bigr)
=
\frac12
\mathbb E\Bigl[
\bigl(h(X_s)-h(Y_s)\bigr)
\bigl(\log X_s-\log Y_s\bigr)
\Bigr]
\ge 0,
\]
where \(Y_s\) is an independent copy of \(X_s\). Hence \(\eta'(s)\ge 0\).
\end{proof}

\begin{lemma}
For every \(r>0\),
\[
-\frac{\zeta'}{\zeta}(1+r)
\le
\frac{\log 2}{2^r-1}.
\]
\end{lemma}

\begin{proof}
For \(s>1\),
\[
\frac{\eta'}{\eta}(s)
=
\frac{\zeta'}{\zeta}(s)
+
\frac{\log 2}{2^{s-1}-1}.
\]
By the previous lemma, \(\eta'(s)\ge 0\), and since \(\eta(s)>0\), we have
\(\eta'(s)/\eta(s)\ge 0\). Therefore
\[
-\frac{\zeta'}{\zeta}(s)
\le
\frac{\log 2}{2^{s-1}-1}.
\]
Taking \(s=1+r\) gives the claim.
\end{proof}

\begin{theorem}
Let \(\Lambda\) denote the von Mangoldt function. For every \(n\ge 1\),
\[
\sum_{q\ge 1}
\frac{\Lambda(q)}{q\log(nq)\log(2nq)}
\le
\frac{1}{\log(2n)}.
\]
Equivalently, since \(\Lambda(1)=0\), the sum may be read as a sum over
\(q\ge 2\).
\end{theorem}

\begin{proof}
For \(x>1\),
\[
\frac{1}{\log x}
=
\int_0^\infty x^{-u}\,du.
\]
Using Tonelli's theorem and the identity
\[
-\frac{\zeta'}{\zeta}(s)
=
\sum_{q\ge 1}\frac{\Lambda(q)}{q^s}
\qquad (s>1),
\]
we obtain
\[
\begin{aligned}
\sum_{q\ge 1}
\frac{\Lambda(q)}{q\log(nq)\log(2nq)}
&=
\int_0^\infty\int_0^\infty
n^{-u-v}2^{-v}
\sum_{q\ge 1}\frac{\Lambda(q)}{q^{1+u+v}}
\,du\,dv  \\
&=
\int_0^\infty\int_0^\infty
n^{-u-v}2^{-v}
\left(-\frac{\zeta'}{\zeta}(1+u+v)\right)
\,du\,dv .
\end{aligned}
\]
Put \(r=u+v\). Then \(0\le v\le r\), and
\[
\int_0^r 2^{-v}\,dv
=
\frac{1-2^{-r}}{\log 2}.
\]
Hence
\[
\sum_{q\ge 1}
\frac{\Lambda(q)}{q\log(nq)\log(2nq)}
=
\frac{1}{\log 2}
\int_0^\infty
n^{-r}(1-2^{-r})
\left(-\frac{\zeta'}{\zeta}(1+r)\right)
\,dr .
\]
By the previous lemma,
\[
-\frac{\zeta'}{\zeta}(1+r)
\le
\frac{\log 2}{2^r-1}.
\]
Therefore
\[
\begin{aligned}
\sum_{q\ge 1}
\frac{\Lambda(q)}{q\log(nq)\log(2nq)}
&\le
\frac{1}{\log 2}
\int_0^\infty
n^{-r}(1-2^{-r})
\frac{\log 2}{2^r-1}
\,dr  \\
&=
\int_0^\infty
n^{-r}\frac{1-2^{-r}}{2^r-1}
\,dr  \\
&=
\int_0^\infty
n^{-r}2^{-r}\,dr  \\
&=
\int_0^\infty
(2n)^{-r}\,dr  \\
&=
\frac{1}{\log(2n)}.
\end{aligned}
\]
This proves the theorem.
\end{proof}

\begin{proposition}
Let \(A\subseteq 2\mathbb N\) be primitive, meaning that no two distinct
elements of \(A\) divide one another. Then
\[
\sum_{a\in A}\frac{1}{a\log a}
\le
\frac{1}{2\log 2}.
\]
Consequently, \(2\) is Erd\H{o}s-strong.
\end{proposition}

\begin{proof}
Let
\[
V:=\{1\}\cup 2\mathbb N.
\]
Define a downward weighted directed graph on \(V\) as follows. First put one
special edge
\[
2\longrightarrow 1
\]
of weight
\[
w(2\to 1):=\frac{1}{2\log 2}.
\]
For every \(n\ge 1\) and every \(q\ge 2\), put an edge
\[
2nq\longrightarrow 2n
\]
of weight
\[
w(2nq\to 2n)
:=
\frac{\Lambda(q)}
{2nq\,\log(nq)\,\log(2nq)}.
\]

For \(x\in V\), let \(O(x)\) be the total weight of the edges leaving \(x\),
and let \(I(x)\) be the total weight of the edges entering \(x\).

We first compute \(O(2n)\). If \(n=1\), the only outgoing edge from \(2\) is
\(2\to 1\), so
\[
O(2)=\frac{1}{2\log 2}.
\]
If \(n>1\), then the outgoing edges from \(2n\) are obtained by writing
\(2n=2(n/d)d\), where \(d\mid n\) and \(d\ge 2\). Therefore, using
\[
\sum_{d\mid n}\Lambda(d)=\log n
\]
and \(\Lambda(1)=0\), we get
\[
\begin{aligned}
O(2n)
&=
\sum_{\substack{d\mid n\\ d\ge 2}}
\frac{\Lambda(d)}
{2n\,\log n\,\log(2n)}  \\
&=
\frac{1}{2n\,\log n\,\log(2n)}
\sum_{d\mid n}\Lambda(d)  \\
&=
\frac{1}{2n\log(2n)}.
\end{aligned}
\]
Thus, for every \(n\ge 1\),
\[
O(2n)=\frac{1}{2n\log(2n)}.
\]

Next compute \(I(2n)\). The incoming edges are
\[
2nq\longrightarrow 2n
\qquad(q\ge 2),
\]
so
\[
\begin{aligned}
I(2n)
&=
\sum_{q\ge 2}
\frac{\Lambda(q)}
{2nq\,\log(nq)\,\log(2nq)}  \\
&=
\frac{1}{2n}
\sum_{q\ge 2}
\frac{\Lambda(q)}
{q\log(nq)\log(2nq)}.
\end{aligned}
\]
By the theorem,
\[
I(2n)
\le
\frac{1}{2n\log(2n)}
=
O(2n).
\]

Now reverse the edges and define a sub-Markov chain. Put
\[
C:=O(2)=\frac{1}{2\log 2}.
\]
The chain starts at \(1\), then moves deterministically from \(1\) to \(2\).
From a vertex \(2n\), it moves to \(2nq\), where \(q\ge 2\), with probability
\[
\mathbb P(2n\to 2nq)
:=
\frac{w(2nq\to 2n)}{O(2n)}.
\]
The total probability of leaving \(2n\) is
\[
\sum_{q\ge 2}\mathbb P(2n\to 2nq)
=
\frac{I(2n)}{O(2n)}
\le 1.
\]
The remaining probability is interpreted as killing the chain.

Let \(P(2n)\) be the probability that the chain ever visits \(2n\). We claim
that
\[
P(2n)=\frac{O(2n)}{C}
\qquad(n\ge 1).
\]
For \(n=1\), this is immediate, because the chain moves from \(1\) to \(2\)
with probability \(1\), and \(O(2)=C\).

Assume \(n>1\), and suppose the identity has been proved for all proper
divisors of \(n\). The immediate predecessors of \(2n\) in the reversed graph
are the vertices \(2n/d\), where \(d\mid n\) and \(d\ge 2\). Hence
\[
\begin{aligned}
P(2n)
&=
\sum_{\substack{d\mid n\\ d\ge 2}}
P(2n/d)\,
\mathbb P(2n/d\to 2n)  \\
&=
\sum_{\substack{d\mid n\\ d\ge 2}}
\frac{O(2n/d)}{C}\,
\frac{w(2n\to 2n/d)}{O(2n/d)}  \\
&=
\frac{1}{C}
\sum_{\substack{d\mid n\\ d\ge 2}}
w(2n\to 2n/d)  \\
&=
\frac{O(2n)}{C}.
\end{aligned}
\]
This proves the claim by induction.

Let \(B\subseteq A\) be finite. Along every non-killed path of the reversed
chain, the current integer strictly increases by divisibility:
\[
2n\longmapsto 2nq
\qquad(q\ge 2).
\]
Since \(A\) is primitive, such a path can visit at most one element of \(A\).
Therefore the events
\[
\{\text{the chain visits } b\},
\qquad b\in B,
\]
are pairwise disjoint. Using the formula for \(P(2n)\), we get
\[
\sum_{b\in B}\frac{O(b)}{C}
=
\sum_{b\in B}P(b)
\le 1.
\]
Thus
\[
\sum_{b\in B}O(b)\le C.
\]
Since \(b\in 2\mathbb N\), the outflow formula gives
\[
O(b)=\frac{1}{b\log b}.
\]
Hence
\[
\sum_{b\in B}\frac{1}{b\log b}
\le
\frac{1}{2\log 2}.
\]
Taking the supremum over all finite subsets \(B\subseteq A\), equivalently
applying monotone convergence, yields
\[
\sum_{a\in A}\frac{1}{a\log a}
\le
\frac{1}{2\log 2}.
\]
Equality is attained by the primitive set \(A=\{2\}\). Therefore \(2\) is
Erd\H{o}s-strong.
\end{proof}

\end{document}
